\subsection*{一步一步建立直观}

\paragraph{1. 把一对实数当作一个可相乘的对象}

复数 $a+bi$ 对应平面中的点，但关键增加的是乘法规则。乘以 $i$ 把实轴正方向转到虚轴正方向，继续乘一次就到负实轴。右图先停在四分之一圈，用位置而非符号记忆来理解这个动作。



\paragraph{2. 两次直角旋转解释平方为负一}

半圈对应两次乘以 $i$，因此 $i^2=-1$ 与平面旋转相容。一般复数乘法则把模长相乘、辐角相加。这种结构让复数适合同时记录缩放和旋转；它不是在实数轴上寻找一个平方为负数的普通实数。



\paragraph{3. 连续的缩放与转动可以合成}

在螺线情形中，每一小段运动同时改变长度和方向。用复数相乘，两个动作可以用一个数统一表示。零复数没有确定辐角，取辐角还存在整圈的不唯一性；这些边界在研究复对数时会再次出现，不能被圆上的直观省略。



\subsection*{定义、命题与证明}

\paragraph{定义：复数的模长}

对 $z=a+bi$，定义 $|z|=\sqrt{a^2+b^2}$。它等于复平面中从原点到该点的距离。

\paragraph{命题：乘以 i 保持长度并旋转直角}

若 $z=a+bi$，则 $iz=-b+ai$，且 $|iz|=|z|$；对非零 $z$，该操作是逆时针直角旋转。

\begin{proof}

由 $i^2=-1$，有 $i(a+bi)=ai-b$。坐标变换为 $(a,b)\mapsto(-b,a)$，新旧长度平方均为 $a^2+b^2$，点积则为 $-ab+ab=0$。

再检查 $(1,0)$ 被送到 $(0,1)$，就确定了旋转方向。零向量同样保持长度，但没有可供比较的辐角。

\end{proof}
